\section{A fixed-aperture first-impact inverse}
\label{sec:aperture}
The resettable apparatus permits a simpler global observation than a list
of protected return laws. A numerical bound on a lattice presentation
places a translate of every obstacle type, and a generating pair of
translations of any one type, in a fixed laboratory square. Short
first-impact launches in that square therefore see a complete periodic
motif. This observation removes the need to normalize a return law by an
unknown cell area. It also removes the matching of independently oriented
contact frames: all impacts already lie in one laboratory frame.

The resulting theorem concerns the same periodic reflecting bodies and
the same collision apparatus. Its datum is a single first-impact
\emph{position law}, not a scalar collision count, a passive orbit, or an
exterior travelling-time spectrum. Its aperture is fixed by the numerical
prior rather than increased with the target accuracy. The local-return
route in the subsequent sections is retained in full; it gives a different
completion certificate when the acquired clouds have not already covered
a periodic motif.

\subsection{A weaker geometric prior and one position law}
\begin{definition}[Translation-separated periodic class]
\label{def:aperture-prior}
Fix numerical bounds
\[
 r_0,D_0,L_0,M_6,\kappa_+<\infty,\qquad
 d_0,\kappa_-,V_0,\sigma>0,\qquad 0<\beta\le1.
\]
Consider nonempty periodic unions of disjoint strictly convex bodies
\[
 \mathcal O=\bigcup_{i=1}^r(C_i+\Lambda),\qquad
 1\le r\le r_0,\quad \Lambda=L\Z^2,
\]
where $\|L\|,\|L^{-1}\|\le L_0$ for some unspecified presentation,
$\covol\Lambda\ge V_0$, and different physical bodies have distance
at least $d_0$. The diameter is at most $D_0$, curvature lies in
$[\kappa_-,\kappa_+]$, and each Steiner-centered support function $p_i$
has $\|p_i\|_{C^{6,\beta}}\le M_6$. The type condition is only
\begin{equation}\label{eq:aperture-type}
                   \|p_i-p_j\|_{L^2(0,2\pi)}\ge\sigma\quad(i\ne j).
\end{equation}
The support functions are compared in the common laboratory directions;
there is no minimization over rotations or reflections in this condition.
No rotation or reflection margin, bridge network, protected return
rectangle, onset bracket, or free-area measurement is prescribed.
\end{definition}

Thus symmetric bodies, including disks, are allowed. Rotated congruent
bodies may be different types provided \eqref{eq:aperture-type} holds.
For one type the condition is empty. Every table in the stronger class
of Definition~\ref{def:prior} belongs to this class with the same or weaker
numerical bounds. The type condition implies that $\Lambda$ is the full
translation group: a translation period must carry any chosen component
to a component with the same centered support; it is therefore one of
that component's declared lattice translates. A compact body has no
nonzero translation symmetry.

Put
\begin{equation}\label{eq:aperture-size}
 B=1+L_0,\qquad \tau=\min(1,d_0/8),\qquad
 R=5B+2D_0+3,\qquad Q_R=[-R,R]^2.
\end{equation}
At every attempt choose a uniform position in $Q_R$ and an independent
uniform direction. A solid start is a failure. From a free start keep
only the first impact before time $\tau$, if one occurs, and otherwise
record failure. Only impact positions in $Q_{5B+D_0+1}$ are retained.
Denote the resulting measure on positions together with failure by
$\mathsf F_{\mathcal O}$. The normalization is by the known square area
$4R^2$, not by the unknown free cell area. For finite data the retained
position has a calibrated enclosure of radius at most $\varepsilon_x$.
No later collision, incoming or outgoing angle, boundary parameter,
normal, travel time beyond the cutoff, or obstacle label is needed.
The common laboratory coordinates and resettable launches remain inputs.

\begin{theorem}[Fixed-aperture rigidity and finite acquisition]
\label{thm:aperture-main}
On Definition~\ref{def:aperture-prior}'s class, the single position law
$\mathsf F_{\mathcal O}$ determines the entire periodic obstacle union,
its full translation lattice, the number of types, and its free cell area.
This is an absolute laboratory-frame determination; forgetting that frame
leaves one common Euclidean isometry, type permutation, representative
changes and an integer change of lattice basis.

There are constants $C,c,\nu_0>0$, depending only on the displayed
numerical prior, with the following finite-data version. For
$0<\nu<\nu_0$ and $0<\delta<1/4$, at most
\begin{equation}\label{eq:aperture-cost}
 C\nu^{-3/4}\log\frac{C}{\nu\delta}
\end{equation}
independent attempts from this same square, with
$\varepsilon_x\le c\nu^{3/2}$, suffice with probability at least
$1-\delta$ to recover the correct number of types and the exact rational
relations of a generating period list. The reconstructed supports and
a period basis have error at most $C\nu$ in the finite-presentation
metric defined below, and the free-area estimate has error at most $C\nu$.
Every failed or discarded launch is charged in \eqref{eq:aperture-cost}.
The method is finite: it uses clustering, hulls, smoothing and finite
rational arithmetic, not an optimization over a class of unknown tables.

Neither the aperture nor the time horizon grows with accuracy. No
auxiliary recentering, return selection, endpoint histogram, free-area
normalization, quantitative asymmetry, or clear-bridge prior is used.
The bound counts launch attempts only; localization precision, arithmetic,
travel and apparatus setup are separate resources. No minimax or passive
trajectory assertion is included.
\end{theorem}

The comparison metric uses one representative support function for each
type in its laboratory placement and one basis matrix for the lattice.
It is the infimum of the maximum $C^2$ support error and matrix error
over one common isometry, type permutations, representative translates,
and integer basis changes. It is not the global Hausdorff distance between
two infinite periodic sets with different period matrices. The proof
produces particular bounded representatives and matched bases realizing
the asserted error; a continuously chosen reduced lattice basis is neither
assumed nor needed.

\subsection{All relevant bodies from one cloud}
A first-hit measure is singular relative to plane area. It must not be
estimated by a two-dimensional smooth density. What matters here is its
mass on each short boundary arc.

For any body meeting $Q_{5B}$, its entire boundary lies in
$Q_{5B+D_0}$, and its exterior collar of width $\tau/2$ lies in $Q_R$.
For a boundary arc $I$, parameterize launches by
\[
 q=x-rv,\qquad \tau/4\le r\le\tau/2,\qquad
                 -\pi/3\le\phi\le\pi/3,
\]
where $v$ makes angle $\phi$ with the inward normal at $x\in I$.
The supporting half-plane excludes an earlier hit on that body, and
$r<d_0$ excludes another body. The Jacobian of the first-hit coordinates
is $\cos\phi$. Consequently, with $S=4R^2$,
\begin{equation}\label{eq:aperture-flux}
       \mathsf F_{\mathcal O}(I)\ge\frac{\tau|I|}{8\pi S}.
\end{equation}
This is an unconditional lower bound per attempted launch. It does not
assume a uniform distribution on the unknown boundary.

There are at most
\[
 M=\left\lceil64\left(1+\frac{R+D_0+d_0}{d_0}\right)^2\right\rceil
\]
bodies meeting the relevant square, by packing disjoint $d_0/2$ disks
about points in different bodies. Their perimeters are at most $\pi D_0$
and at least $2\pi/\kappa_+$. Choose
\[
 0<q<\min(d_0/40,1,\pi/\kappa_+),\qquad
 \varepsilon_x<\min(d_0/100,q/10,1/16).
\]
Partition each boundary into arcs of lengths between $q/2$ and $q$.
From \eqref{eq:aperture-flux} a coupon bound gives simultaneous coverage
with failure probability at most $\delta$ as soon as
\begin{equation}\label{eq:aperture-coupon}
 N\ge\frac{16\pi S}{\tau q}
       \log\frac{2M(1+\pi D_0/q)}{\delta}.
\end{equation}
Connect reported points at distance less than $d_0/3$. On this event,
no component joins different bodies, and each relevant body has a
connected complete cloud. Consecutive sampled boundary points have
arclength gaps at most $2q$, whereas different bodies are separated by
$d_0-2\varepsilon_x$. If $P$ is a complete cloud hull, then
\begin{equation}\label{eq:aperture-hull}
 d_H(P,C)=\|p_P-p_C\|_\infty
                 \le e:=\kappa_+q^2/2+\varepsilon_x.
\end{equation}
At a support point the first tangential derivative of the support height
vanishes and the second derivative in arclength is bounded by
$\kappa_+$, which proves this inequality. Localization adds at most
$\varepsilon_x$ to every support value. This is the classical boundary
hull mechanism, here supplied by a physical first-hit minorization;
compare \cite{Schneider,PSSW}.

Only complete clouds may be used for shape recognition. The following
selection avoids assuming that completeness is directly observable.
Write $s(P)$ for the Steiner point of any cluster hull, and keep the
clusters satisfying $\|s(P)\|_\infty\le4B$. Each hull lies in the
$\varepsilon_x$-neighborhood of its true convex body, even for a partial
outer cloud. Since $s(P)\in P$, a kept cluster's body meets
$Q_{4B+\varepsilon_x}\subset Q_{5B}$. The simultaneous event above then
ensures complete coverage and connectivity of that very body. Thus
partial outer components cannot masquerade as additional inner types.
This argument uses raw hulls, not signed smooth approximations, in the
selection test.

\subsection{Recovering the full lattice before any area calculation}
The bounded-presentation hypothesis provides a coverage theorem, not
numerical period labels. For every type, some translate of its Steiner
point lies in $L[-1/2,1/2]^2$, hence has Euclidean norm at most
$L_0/\sqrt2<B$. Since $|s(P)-s(C)|\le2e$, the selected clouds contain
at least one copy of every type when $2e<1/8$.

Temporarily suppose these selected bodies and their centered supports
are known with error at most $u$ in $C^2$, and their centers with error
at most $t$. When $2\sqrt{2\pi}u<\sigma/4$, the threshold $\sigma/2$
on centered $L^2$ support differences separates their exact translation
types: copies agree in the same laboratory directions, while different
types satisfy \eqref{eq:aperture-type}. No orthogonal alignment is
performed. In particular a disk does not cause an orientation singularity.

Choose any cluster with raw center in $Q_{2B}$ as a root. Such a cluster
exists by the preceding coverage argument. Take differences of centers
between all selected copies of its type and the root. Call the true
vectors $d_j$ and their approximations $\widehat d_j$. Every $d_j$ is
a true period. More importantly, if $\ell_1,\ell_2$ are the columns of
the bounded presentation, the root's copies translated by $\ell_i$ have
true centers in $Q_{3B+1/8}$ and estimated centers in $Q_{3B+1/4}$.
They are therefore selected within $Q_{4B}$. Their differences are
$\ell_1,\ell_2$. It follows that
\begin{equation}\label{eq:aperture-generation}
                         \sum_j\Z d_j=\Lambda.
\end{equation}
Any cutoff ambiguity concerns optional additional copies only; it cannot
remove these protected generators or introduce a nonperiod. The bounded
aperture thus certifies saturation and absence of unseen types directly.
There is no area-defect test in this route.

All these differences have norm at most $U=12B$ for small errors.
Set
\begin{equation}\label{eq:aperture-Q}
                     Q=\max\{1,\lceil U^2/V_0\rceil\}.
\end{equation}
Every nonzero determinant of two true periods has magnitude at least
$V_0$, and its perturbation is at most $(2U+t')t'$ when vector errors
are at most $t'$. A threshold $V_0/2$ therefore detects independence
for sufficiently small $t'$. Select a pair with maximal estimated
absolute determinant and write $D=[d_1\ d_2]$. Its index in $\Lambda$
is an integer
\[
                    n=|\det D|/\covol\Lambda\le Q.
\]
Each coordinate of $D^{-1}d_j$ is a bounded rational whose reduced
denominator divides $n$. Distinct rationals with denominator at most
$Q$ differ by at least $Q^{-2}$. Inversion of this $2\times2$ matrix,
using $|\det D|\ge V_0$ and the bound $U$, gives a computable $C_*$
such that coordinate errors are at most $C_*t'$. Once this is less
than $1/(3Q^2)$, finite rational search recovers every coordinate
exactly. One can restrict the numerator to the known coordinate bound
$2U^2/V_0+1$ times $Q$. No integer span of noisy real vectors is taken.

Let $q_*$ be the least common denominator of the recovered coordinates;
it divides $n$ and is at most $Q$. A column Hermite basis $H$ of the
integer group generated by $q_*I_2$ and the vectors $q_*D^{-1}d_j$
then gives
\begin{equation}\label{eq:aperture-lattice}
          \Lambda=(D H/q_*)\Z^2,\qquad
          \widehat L=\widehat D H/q_*.
\end{equation}
The group contains $q_*\Z^2$, so its Hermite determinant is at most
$q_*^2$. All entries of $H$ are bounded by $Q^2$ in the usual
nonnegative remainder convention. This is a finite family, and the
period-basis error is consequently at most $C t'$. Optional repetitions,
different choices of root or independent pair, and cutoff changes only
choose another basis of the same exact lattice. The metric permits
that basis change rather than asserting that a reduced basis varies
continuously.

Keep one body per recovered type. These bodies and
\eqref{eq:aperture-lattice} determine the entire table. The cell's free
area can now be estimated geometrically by
\begin{equation}\label{eq:aperture-area}
              \widehat A=|\det\widehat L|-\sum_i\area(P_i),
\end{equation}
where $P_i$ are the selected raw hulls. On bounded convex bodies the
area error is $O(e)$ under a Hausdorff perturbation $e$, by the parallel
body formula and the uniform perimeter bound. Since the center errors
are $O(e)$, the matrix and free-area errors here are in fact $O(e)$
after type and rational recovery. The $C\nu$ conclusion of the theorem
is a common bound for these errors and the smoother support error.
The scalar estimate \eqref{eq:aperture-area} is not defined by the area
of a separately smoothed numerical curve.

\subsection{Two-derivative confidence and proof of the theorem}
For clarity, we state the smoothing step independently of the return-law
construction. Use the even nonnegative kernel
\[
 \varphi(z)=\frac{315}{256}(1-z^2)^4\one_{[-1,1]}(z),\qquad
 K_h=\tfrac43\varphi_h-\tfrac13\varphi_{2h}.
\]
Its mass is one and its moments of orders one, two and three vanish.
Write $C_\varphi=1+\|\varphi'\|_1+\|\varphi''\|_1\le512$.
For a periodic support function with six derivatives bounded by $B_6$,
\begin{equation}\label{eq:aperture-smoothing}
       \|K_h*p_P-p_C\|_{C^2}
                    \le2C_\varphi e h^{-2}+B_6h^4.
\end{equation}
Differentiating the kernel gives the first term. Taylor expansion of
$p_C^{(j)}$ through degree three, for $j\le2$, gives the second,
since the absolute fourth moment is at most $(20/3)h^4$ and the
Taylor denominator is $24$. These calculations do not differentiate
the noisy polygon. Although $K_h$ is signed, an error less than
$1/(4\kappa_+)$ implies
\[
       K_h*p_P+(K_h*p_P)''>1/(2\kappa_+),
\]
so the reconstructed support still defines a regular strictly convex
body. Convexity is justified by this inequality, not by positivity of
the convolution weights.

Translate by the raw Steiner point before applying this estimate, and
restore the same point afterwards. A uniform $B_6$ follows from $M_6$
and the bounded centering error. The linear Steiner-point functional is
bounded in the support sup norm, so recentering for type comparisons
changes errors by at most a uniform factor. Set, for small $\nu$,
\[
 h=(\nu/(8B_6))^{1/4},\quad
 E_\nu=\frac{\nu h^2}{128C_\varphi},\quad
 q=\sqrt{E_\nu/(4\kappa_+)},\quad
 \varepsilon_x\le E_\nu/8.
\]
Decrease $\nu_0$ to enforce all coverage, selection, convexity, type and
rational thresholds already stated. Then $e\le E_\nu/4$ and
\eqref{eq:aperture-smoothing} is less than $\nu/4$. Deterministic
quadrature can be performed to a further fixed fraction of $\nu$:
the hull has finitely many known linear pieces in its support, the
kernel is explicit, and the integrals have computable bounds. The raw
hull's area and Steiner point can likewise be bounded to a smaller
fraction of $e$. This gives a finite numerical construction with all
arithmetic errors budgeted.

Here $q$ is a constant times $\nu^{3/4}$ and $E_\nu$ a constant times
$\nu^{3/2}$. Substitution into \eqref{eq:aperture-coupon} proves
\eqref{eq:aperture-cost}. On its one simultaneous coverage event, all
subsequent classification, generating-list and rational tests are
correct. No additional probability is spent on data-dependent root or
basis choices: the deterministic argument holds for every choice from
the same covered cloud. The reconstructed finite presentation has
positive curvature and separation for small $\nu$. Only bounded integer
translates can approach a bounded fundamental region, since the matched
lattice basis and its inverse are bounded by the prior and the finite
Hermite family; separation checking therefore has no missing distant
translate. This proves the finite-data assertion.

For exact laws, \eqref{eq:aperture-flux} shows that the support in the
inner region is exactly the union of the complete relevant boundaries.
There are no other impacts, and separation identifies its connected
boundary components and their bounded interiors. Perform the preceding
selection, type identification and period generation without errors.
Every type and a lattice basis are present by the coverage argument.
Equations \eqref{eq:aperture-generation} and
\eqref{eq:aperture-area} then give the exact periodic union and its free
area. This proves the rigidity assertion of
Theorem~\ref{thm:aperture-main}.

\begin{corollary}[Simultaneous certificates without moving the aperture]
\label{cor:aperture-anytime}
Let $\nu_j=2^{-j}\nu_0/2$ and spend
$\delta_j=\delta/[2(j+1)^2]$ on fresh batches from the same square,
using the preceding localization and launch bounds. With probability
at least $1-\delta$ every issued finite-presentation certificate is
correct. By the time target $\nu_j$ is reached, the cumulative number
of attempts is at most $C\nu_j^{-3/4}\log(C/(\nu_j\delta))$.
\end{corollary}
\begin{proof}
The sum of the failure allowances is less than $\delta$ for $j\ge1$.
The simultaneous assertion follows by a union bound. The budgets grow
geometrically, and their logarithmic factors grow at most linearly in
$j$ plus $\log(1/\delta)$; summing gives the displayed bound. Each
batch uses fresh positions with its stated localization, rather than
pretending to refine a previously recorded coarse coordinate. A finite
certificate may be withheld when numerical tests are indeterminate;
on the coverage event the declared thresholds ensure they pass.
\end{proof}

\subsection{Scope and the retained intrinsic law}
For a single disk repeated on a lattice, every rotation and reflection
is a body symmetry, yet the theorem recovers its center translates and
full lattice. Several disks of distinct radii also satisfy
\eqref{eq:aperture-type} whenever the numerical bounds permit their
separated placement. These examples are outside the quantitative
asymmetry assumption of the record-frame theorem. They demonstrate a
strict extension of its geometric scope, not a reformulation of an
asymmetric contact descriptor.

The bounded-presentation hypothesis and translation-type gap are still
substantive. An arbitrarily small local field of view cannot enumerate
unknown periodic types. Nor can arbitrary identical motifs in a
nonprimitive declared cell be assigned an observable cell multiplicity
without further information. We recover the full translation group of
the physical union in the stated type-separated class, not an externally
chosen multiple cell. The position law also retains a common laboratory
frame; it does not solve passive, count-only or marked-length rigidity.

The experiment's reduction is precise: first impacts in one fixed square
are a restriction of the already allowed resettable collision apparatus.
No hypothetical boundary, distance or normalized phase-measure sensor
has been added. The new proof replaces return-law completion by certified
coverage using the existing numerical lattice bound. The intrinsic
identity, local inverse, bridge fingerprint, area-defect certificate,
sequential revalidation and both cloud-cost theorems remain active below.
Their use of the area normalization is indispensable for their
\emph{record-local} information set, not for a globally covering aperture.
The established convex-hull and smoothing rates are ingredients, not a
claim of a new general random-polytope theorem; the conclusion here is
the finite first-impact determination of the complete periodic geometry.
